% Candidate report section — Track K (2026-11-22).
% Source: reports/K02_physics_model_refinement.md and results/k02_*.
% Numbers that must be regenerated before this file is moved into report/:
%   @M4RMS@ @M4DELTA@ @M4MEAN@ @SELECTED@ @V3THR@ @V3FR@ @THRESHDELTA@ @DETECTION
% The Track K reports are the source of truth; this section is a candidate, not merged.

\section{Offline airframe-model refinement and the calibration mode}
\label{sec:k-physics}

The published clip floor (14.73 deg) is the \emph{drag-free} mode of the required lean;
re-running A2's own construction reproduces it to the printed precision, together with the
other three modes (8.96, 12.95, 6.46 deg) and their detection floors.  Schema v2's
\texttt{lean\_req} is the drag-free lean to a median 0.016 deg, and the consumer form
\texttt{lean\_req}$+k_d$\texttt{vel\_req} is a linearisation that differs from the exact
flatness axis by a median 2.2 deg and a p95 20.4 deg on testable frames.  This resolves the
published label conflict: the drag-free floor is correctly labelled, and the ``drag-free''
anchor that appeared to give 31.1 deg is the $k_d=-0.295$ arm of a shifted reparameterisation.

We then asked whether an offline model of how the airframe realises a commanded
acceleration lowers the floor.  Five candidates were pre-registered in order of complexity:
per-axis linear drag (M1), M1 plus a first-order attitude lag (M2), M1 plus quadratic drag
(M3), body-frame rotor drag (M4, after Faessler et al.\ \cite{faessler2018flatness}), and a
strongly regularised ridge residual on $(a,v,$jerk$)$ fitted offline and fixed at test time.
Every model was fitted on path kinematics and truth attitude only, leave-one-flight-out,
and every parameter was fixed before any clipped edit was scored.

\begin{table}[t]
\centering
\caption{Leave-one-flight-out angular RMS between the actual thrust axis and the model
required axis, eight TII flights.  The frozen per-flight lag nuisance is given to every
model.  Frame-weighted pooled RMS; per-flight values are in the Track K report.}
\label{tab:k-models}
\begin{tabular}{lccc}
\toprule
model & pooled RMS (deg) & change vs M0 & fitted parameters \\
\midrule
M0 current            & 3.648 & ---      & $k_d=0.295$ \\
M1 per-axis linear    & 3.612 & $-1.0\%$ & $k_f=0.27$, $k_l=0.23$ \\
M2 M1 + attitude lag  & 3.573 & $-2.1\%$ & $\tau=+12.5$ ms \\
M3 M1 + quadratic     & 3.665 & $+0.5\%$ & $k_q\approx 4$--$9\times10^{-4}$ \\
M4 rotor drag         & @M4RMS@ & @M4DELTA@ & @M4PARAMS@ \\
M5 ridge residual     & 7.40  & $+103\%$ & $\alpha=0.1$ (inner LOFO) \\
\bottomrule
\end{tabular}
\end{table}

No candidate approaches the pre-registered 25\,\% RMS reduction, and the best
(@SELECTED@) leaves the gate-kept-dev F0 threshold at @V3THR@ deg against 8.87 deg for the
exact $k_d=0.295$ required side: a change of @THRESHDELTA@ deg, short of the 1.5 deg
criterion.  The learned residual does not generalise: its selected regularisation more than
doubles the held-out RMS, and at $\alpha=100$ it collapses to M1, because the held-out
racing flights lie outside its feature range.

The reason is measured rather than argued.  The M0 residual has 40--74\,\% of its variance
in 0.5--2 Hz and 17--52\,\% in 2--10 Hz, less than 1\,\% above 10 Hz on TII, and its
magnitude correlates with jerk (0.48--0.77) and speed on the racing flights.  That is a
dynamic attitude response, but none of per-axis, quadratic or body-frame drag reproduces it,
and a first-order lag or lead does not either.  The practical conclusion is that the
existing drag term \emph{is} the physics refinement that the drag-free floor was missing:
the exact $k_d=0.295$ F0 threshold on the gate-kept dev clips is 8.87 deg, below the
handoff's ``all of the mismatch removed'' bracket of 9--9.3 deg, so the remaining floor is
not a drag-model mismatch.

Schema v3 stores the exact model required lean \emph{and} the required thrust axis at the
clip's frame times, with explicit semantics (no $(k_d-0.295)$ rescaling), in new files; the
axis is necessary because an inverted required axis cannot be represented by a lean.
@DETECTION

\begin{thebibliography}{9}
\bibitem{faessler2018flatness}
M.~Faessler, A.~Franchi, and D.~Scaramuzza,
``Differential flatness of quadrotor dynamics subject to rotor drag for accurate tracking
of high-speed trajectories,'' \emph{IEEE Robotics and Automation Letters}, 2018.
arXiv:1712.02402.
\end{thebibliography}
